\Askhsh{23215}{1}
\begin{ylh}{1}
    \item [Διαφορικός Λογισμός (Ρολ, ΘΜΤ, Μονοτονία, Ακρότατα)]
\end{ylh}
Δίνεται συνάρτηση $f: \mathbb{R} \to \mathbb{R}$ η οποία είναι παραγωγίσιμη και ισχύει $f'(x) \neq 0$ για κάθε $x \in \mathbb{R}$. Επίσης, η $f'$ είναι συνεχής.

\begin{alist}
\item Δείξτε ότι η συνάρτηση $f$ είναι 1-1. \monades{7}

\item Δίνεται επιπλέον ότι $f(0)=-1$ και $f(2)=1$.
\begin{rlist}
\item Δείξτε ότι υπάρχει τουλάχιστον ένα $\xi_1 \in (0, 2)$ τέτοιο ώστε $f'(\xi_1) = 1$. \monades{7}
\item Να δώσετε μία γεωμετρική ερμηνεία του συμπεράσματος του ερωτήματος β.i. \monades{4}
\end{rlist}

\item
\begin{rlist}
\item Με βάση τα δεδομένα των προηγούμενων ερωτημάτων, δείξτε ότι η $f$ είναι γνησίως αύξουσα στο $\mathbb{R}$. \monades{5}
\item Αν επιπλέον δίνεται ότι $f'(0) \cdot f'(2)=-1<0$ και η $f'$ είναι γνησίως αύξουσα, δείξτε ότι η εξίσωση $f'(x)=0$ έχει μία μοναδική ρίζα $x_o \in (0,2)$. \monades{9}
\item Έστω $g(x) = f'(x)$. Μελετήστε την συνάρτηση $g$ ως προς τη μονοτονία και τα ακρότατα και συγκεντρώστε τα αποτελέσματα σε έναν πίνακα. \monades{12}
\end{rlist}
\end{alist}